


\begin{frame}{Why Ambiguous!}
   \only<1>{

   \frametitle{Principles of CODA}
      \begin{itemize}
       \item 
   \textbf{scale invariance},
   \item
   \textbf{permutation invariance}, and 
   \item \textbf{subcompositional coherence}
   If \( \Delta_p(\mathbf{x}, \mathbf{y}) \) denotes a distance measure between compositions with \( p \) parts, then
    \begin{equation*}
    \Delta_k(\mathbf{x}, \mathbf{y}) \geq \Delta_d(\mathbf{x}_d, \mathbf{y}_d),
    \end{equation*}
    where \( \mathbf{x}, \mathbf{y} \) are compositions with \( k \) parts, and \( \mathbf{x}_d, \mathbf{y}_d \) are corresponding subcompositions with \( d < k \) parts.
    \end{itemize}}
\only<2>{
\textbf{\textcolor{blue}{Question 1:}} Are these principles designed to achieve consistent results with the original data $\mathbf{Y}$?  

The arguments supporting $\mathbf{a}(\mathbf{Y})$ are rarely evaluated for consistency with the original vector $\mathbf{Y}$, even though Aitchison himself criticized standard statistical approaches for inconsistency with the original data.
}

\only<3>{
As an illustration, consider
\[
\mathbf{S}(\mathbf{Y}) = \Big( \frac{Y_1}{Y_1 + Y_2}, \frac{Y_2}{Y_1 + Y_2} \Big) = (U, 1-U).
\]
\textbf{\textcolor{blue}{Can we assume that \(U\) follows a Beta distribution?}}

\textbf{\textcolor{red}{Aitchison's critique:}} For \(U\) to follow a Beta distribution,
\[
Y_i \sim^{\text{iid}} \text{Gamma}(\alpha_i, \lambda), \quad i=1,2.
\]
This generalizes naturally to the Dirichlet distribution.

His proposal instead:
\[
\left(\frac{1}{2}\ln\frac{Y_1}{Y_2}, \frac{1}{2}\ln\frac{Y_2}{Y_1}\right)
\]
}

\only<4>{
\textbf{\textcolor{blue}{Question 2:}} Are these principles designed to provide non-spurious correlations?

Correlations among components of $\mathbf{a}(\mathbf{Y})$ are themselves spurious.
}

\only<5>{
\textbf{Spurious correlation example:}

Even if  
\[
\operatorname{corr}(Y_1, Y_2)=0,
\]
we usually observe
\[
\boxed{
\operatorname{corr}\Big(\frac{Y_1}{S}, \frac{Y_2}{S}\Big) \neq 0, \quad S=\sum_{j=1}^{D} Y_j
}
\]
\textbf{Reason:} shared denominator induces dependency.
}

\only<6>{
\[
\boxed{
\operatorname{corr}\Big(\frac{Y_1}{S}, \frac{Y_2}{S}\Big)
=
\frac{
\mathbb{E}\Big(\frac{Y_1 Y_2}{S^2}\Big) - \mathbb{E}\Big(\frac{Y_1}{S}\Big)\mathbb{E}\Big(\frac{Y_2}{S}\Big)
}{
\sqrt{\operatorname{Var}\Big(\frac{Y_1}{S}\Big)} \sqrt{\operatorname{Var}\Big(\frac{Y_2}{S}\Big)}
}
}
\]

\textbf{Key point:} The formula is mathematically valid; the issue is interpretation:
\[
\operatorname{corr}\Big(\frac{Y_1}{S}, \frac{Y_2}{S}\Big) \not\Rightarrow \operatorname{corr}(Y_1,Y_2)
\]
}

\only<7>{
\textbf{Frequently emphasized in CODA:}

\begin{block}{Different denominators}
\[
\left(\frac{Y_1}{S_4},\frac{Y_2}{S_4},\frac{Y_3}{S_4},\frac{Y_4}{S_4}\right),\quad
\left(\frac{Y_1}{S_3},\frac{Y_2}{S_3},\frac{Y_3}{S_3}\right)
\]
\end{block}

\begin{alertblock}{Claimed problem}
\[
\operatorname{corr}\Big(\frac{Y_1}{S_4}, \frac{Y_2}{S_4}\Big) \neq \operatorname{corr}\Big(\frac{Y_1}{S_3}, \frac{Y_2}{S_3}\Big)
\]

\textbf{Both correlations are spurious.}  
It is unclear why equality between two spurious correlations should matter.
\end{alertblock}
}

\only<8>{
\begin{block}{ILR (or Other) Transformed Coordinates}

\begin{minipage}{0.48\textwidth}
\centering
\textbf{} \\[2mm]
\scalebox{0.75}{$
\mathbf{T}_4 = T\Big(\frac{Y_1}{S_4},\frac{Y_2}{S_4},\frac{Y_3}{S_4},\frac{Y_4}{S_4}\Big):(T_1,T_2,T_3,T_4)
$} \\[2mm]
\scalebox{0.75}{$
\mathbf{T}_3 = T\Big(\frac{Y_1}{S_4},\frac{Y_2}{S_4},\frac{Y_3}{S_4}\Big):(T_1,T_2,T_3)
$}
\end{minipage}
\hfill
\begin{minipage}{0.48\textwidth}
\centering
\textbf{} \\[2mm]
\scalebox{0.75}{$
\mathbf{T}'_4 = T\Big(\frac{Y'_1}{S'_4},\frac{Y'_2}{S'_4},\frac{Y'_3}{S'_4},\frac{Y'_4}{S'_4}\Big):(T'_1,T'_2,T'_3,T'_4)
$} \\[2mm]
\scalebox{0.75}{$
\mathbf{T}'_3 = T\Big(\frac{Y'_1}{S'_4},\frac{Y'_2}{S'_4},\frac{Y'_3}{S'_4}\Big):(T'_1,T'_2,T'_3)
$}
\end{minipage}

\end{block}


\begin{alertblock}{Subcompositional Coherence}
\scalebox{0.85}{\[
D_3(\mathbf{T}_3,\mathbf{T}'_3) \le D_4(\mathbf{T}_4,\mathbf{T}'_4)
\]}
\begin{center}
\textbf{\large Surprisingly, the correlations between components remain \textcolor{red}{spurious}, 
but we cannot work in the simplex space accordingly.}
\end{center}
\end{alertblock}
}

\only<9>{
\textbf{\textcolor{blue}{Question 2:}} Are these principles designed to provide non-spurious correlations?  The answer is no.
}

\only<10>{
\textbf{\textcolor{blue}{Question 3:}} Are these principles designed to satisfy theoretical requirements?

It is important to note that \(\mathbf{s}(\cdot)\)  is subset of real space and therefore inherit its standard geometric structure, rendering all conventional mathematical operations well-defined.
}

\only<11>{
\textbf{\textcolor{blue}{Question 4:}} Are these principles widely accepted?  The answer is no.
}

\end{frame}



\begin{frame}{Critics of CoDA: Early and Ongoing}
A number of authors have questioned the strict log-ratio framework
promoted by the “CoDa School.” Major groups include:
\scriptsize   % shrinks everything but remains readable


\begin{itemize}\itemsep1pt
    \item \textbf{Watson}: objections to distance measures and geometry
    \item \textbf{Zier \& Rehder}: challenges to mathematical definitions
    \item \textbf{Butler \& Glasbey}: alternative approaches for zeros
    \item \textbf{Stewart \& Field}: practical issues in real applications
    \item \textbf{Tsagris et al.}: limits for highly variable data
    \item \textbf{Scealy \& Welsh}: critique of the Principles themselves
\end{itemize}



\end{frame}




\begin{frame}{Key Critics: D.F. Watson}
\scriptsize   % shrinks everything but remains readable

    \begin{itemize}
        \item \textbf{D.F. Watson} was a central figure in early debates against Aitchison's log-ratio methods.
        \item **Main Criticism Focus:**
        \begin{itemize}
            \item Questioned the fundamental need for "scale invariance," a core principle of the CoDa School.
            \item His exchanges with Aitchison in \textit{Mathematical Geology} (1989-1992) were instrumental in the development of the CoDa Principles, though from a critical stance.
            \item The text notes that "Watson \& Philip (1989) ... did not see the need for scale invariance."
        \end{itemize}
        \item Tragically, his contributions were cut short due to a contentious exchange where Aitchison used derogatory language.
    \end{itemize}
\end{frame}
\begin{frame}{Key Critics: Zier \& Rehder}
\scriptsize   % shrinks everything but remains readable

    \begin{itemize}
        \item \textbf{Zier \& Rehder} (1998, 2001) continued to raise points against the CoDa School's established views.
        \item **Main Criticism Focus:**
        \begin{itemize}
            \item Challenged technical issues, particularly related to the mathematical definition of a metric within compositional data analysis.
            \item Like Watson, they "did not see the need for scale invariance," arguing it problematic if compositions are strictly defined as summing to one.
            \item Their arguments highlighted that even Aitchison's own earlier work might not have strictly adhered to the current CoDa definition of compositions and scale invariance.
        \end{itemize}
        \item Despite making "some good points," their criticisms were "dismissed relatively easily" by the CoDa School, partly due to perceived technical mistakes on their part.
    \end{itemize}
\end{frame}
\begin{frame}{Key Critics: Butler \& Glasbey}
\scriptsize   % shrinks everything but remains readable

    \begin{itemize}
        \item \textbf{Butler \& Glasbey} (2008) represent a more recent wave of critics whose methods were deemed "unapproved" by the CoDa School.
        \item **Main Criticism Focus:**
        \begin{itemize}
            \item Developed "otherwise sensible analysis" methods that did not conform to the CoDa Principles.
            \item Were compelled to "list and explain their 'failings' relative to the Principles."
            \item Their work highlighted that theoretical difficulties with the Principles meant their methods could only be used as "diagnostic tool or for exploratory analysis" according to CoDa proponents.
            \item This illustrates how the Principles act as a barrier to new methodological developments.
        \end{itemize}
    \end{itemize}
\end{frame}
\begin{frame}{Key Critics: Stewart \& Field, Tsagris et al., Scealy \& Welsh}
\scriptsize   % shrinks everything but remains readable

    \begin{itemize}
        \item \textbf{Stewart \& Field (2010), Tsagris, Preston \& Wood (2011), and Scealy \& Welsh (2011a, 2011b)} also demonstrate the practical limitations of the CoDa School's exclusive approach.
        \item **Main Criticism Focus:**
        \begin{itemize}
            \item Provided "interesting real examples where important practical differences produced by a variety of methods are highlighted."
            \item Showed that the differences between various methods become particularly important for \textbf{highly variable data near or on the boundary of the simplex}.
            \item \textbf{Scealy & Welsh (2011a, 2014, the source text)} specifically critique the **motivation, justification, and relevance of the Principles themselves**, particularly subcompositional coherence.
            \item They argue against the logical basis for scale invariance and subcompositional dominance, proposing alternatives like "amalgamations" as more suitable 'marginals'.
        \end{itemize}
    \end{itemize}
\end{frame}
\begin{frame}{Conclusion on Criticisms}
\scriptsize   % shrinks everything but remains readable

    \begin{itemize}
        \item These critics collectively highlight significant issues:
        \begin{itemize}
            \item The contested definition of a "composition."
            \item The rigid and sometimes illogical foundations of the CoDa Principles (e.g., scale invariance, subcompositional coherence).
            \item The active dismissal of valid alternative methods simply for not conforming to these Principles.
            \item The practical limitations of the CoDa approach, especially with challenging data.
        \end{itemize}
        \item The text emphasizes that the existence and application of these Principles in such a restrictive way is "unusual, perhaps even unique, in statistical data analysis."
    \end{itemize}
\end{frame}



\begin{frame}{The Aitchison-Watson Debate: A Contentious Beginning}
\scriptsize   % shrinks everything but remains readable

    \begin{itemize}
        \item The foundational "Principles of Compositional Data Analysis" largely evolved from a series of intense exchanges.
        \item Primarily, this involved papers and letters between \textbf{J. Aitchison} (as the defender/proponent of log-ratio methods) and \textbf{D.F. Watson} (as a critic of the nascent CoDa School's approach).
        \item These debates, published in \textit{Mathematical Geology} from 1989-1992, set the stage for much of the field's development and controversies.
        \item However, the intellectual exchange took a negative turn:
        \begin{itemize}
            \item Aitchison is noted to have used the word "stupid" in reference to Watson's arguments.
            \item This pejorative remark led to \textbf{Watson discontinuing his contributions} to the debate, effectively silencing a critical voice.
        \end{itemize}
        \item This event highlights how personal attacks can impede scientific discourse and shape the trajectory of a research field.
        \item Scealy & Welsh (2014) emphasize that "the exchange between Watson and Aitchison had a negative impact on both sides and contributed directly to the situation in modelling and analysing compositional data that we have today."
    \end{itemize}
\end{frame}
